% !TeX root=tukedip.tex
\section{Formulácia úlohy}

Zadanie diplomovej práce definuje päť úloh. Táto kapitola rozvádza spôsob, akým sú jednotlivé úlohy riešené, a~uvádza podmienky a~obmedzenia riešenia. Prepojenie úloh s~kapitolami práce sumarizuje tabuľka~\ref{tab:tasks_mapping}.

\subsection{Úloha~1: Analyzovať aktuálny stav moderných prístupových systémov}

Analýza aktuálneho stavu pokrýva existujúce riešenia od jednoduchých systémov na báze UID a~protokolu Wiegand, cez otvorený štandard OSDP v2 s~AES-128 Secure Channel, až po priemyselné systémy s~proprietárnymi protokolmi. Pre každú kategóriu sa identifikujú bezpečnostné vlastnosti a~nedostatky relevantné pre kontext tejto práce. Ako srovnávacia základňa pre experimentálne overenie slúži štandardné nasadenie AEAD podľa RFC~8439 (ChaCha20-Poly1305), ktoré definuje algoritmus a~požiadavky na nonce, no nepredpisuje spôsob generovania nonce, deriváciu kľúčov pre viacero zariadení ani runtime detekciu anomálií.

\subsection{Úloha~2: Analyzovať problematiku bezpečnosti šifrovacích algoritmov}

Teoretická analýza sa zameriava na kryptografické primitíva priamo použité v~prototype: AEAD konštrukcie ChaCha20-Poly1305 a~XChaCha20-Poly1305, ich bezpečnostné modely (IND-CPA, INT-CTXT, nonce-misuse resistance), HMAC-SHA-256 pre autentifikáciu, HKDF pre deriváciu kľúčov a~konštrukciu tamper-evident auditného logu. Porovnanie s~AES-GCM zdôvodňuje voľbu algoritmov pre embedded platformu bez hardvérovej akcelerácie. Analýza identifikuje tri stratégie generovania nonce (náhodný, deterministický HMAC, deterministický s~runtime verifikáciou) ako hlavnú experimentálnu os.

\subsection{Úloha~3: Navrhnúť model prístupového systému}

Model je realizovaný ako funkčný prototyp: hardvérová čítačka (ESP32 s~modulom RC522) komunikuje so serverovou aplikáciou v~jazyku C++20 prostredníctvom HTTP požiadaviek autentifikovaných HMAC-SHA-256. Server implementuje viacvrstvovú ochranu: interný AEAD frame s~AAD binding, HKDF per-reader deriváciu kľúčov, anti-replay sliding window, detektor protokolových anomálií s~karanténou a~tamper-evident audit log. Architektúra oddeľuje spracovanie požiadaviek od rozhodovacej logiky prostredníctvom HTTP rozhrania medzi dvoma nezávislými službami (frontend a~DecisionService), čo umožňuje nezávislé testovanie bezpečnostných vlastností AEAD vrstvy. Šesť konfiguračných profilov (dva algoritmy krát tri nonce stratégie) vytvára experimentálnu maticu, kde všetky ostatné mechanizmy zostávajú konštantné.

\subsection{Úloha~4: Vyhodnotiť účinnosť algoritmov a~navrhnúť modifikácie}

Vyhodnotenie prebieha experimentálne pomocou C++ harnessu, ktorý zdieľa produkčný kód a~umožňuje reprodukovateľné merania. Sedem scenárov testuje odolnosť voči replay útoku, manipulácii hlavičky, nonce enforcement, cross-reader izolácii kľúčov, rollbacku sekvenčného čísla, verifikácii AEAD tagu a~nonce tamper útoku na kanáli. Modifikácia oproti štandardnému nasadeniu RFC~8439 spočíva v~deterministickej konštrukcii nonce (HMAC z~kontextu hlavičky), runtime verifikácii nonce na serveri a~automatickej karanténe zariadení pri detekcii anomálie. Experimenty bežia v~dvoch režimoch: in-process (čisté kryptografické metriky) a~E2E cez HTTP endpoint (validácia celej spracovateľskej cesty).

\subsection{Úloha~5: Vypracovať dokumentáciu podľa štandardov KPI}

Práca je vypracovaná podľa šablóny a~štandardov KPI FEI TUKE. Zdrojový kód, experimentálne dáta a~analytický pipeline sú súčasťou priloženého média.

\subsection{Podmienky a~obmedzenia riešenia}

Prototyp vedome zjednodušuje niektoré aspekty produkčného nasadenia. Karta pracuje v~jednoduchom režime --- UID vystupuje ako identifikátor, nie ako kryptografický dôkaz identity. Úložisko je postavené na SQLite, čo je praktické pre prototyp a~experimenty, no neposkytuje produkčné škálovanie. Povinné TLS nie je v~základnom režime nasadené; dôvernosť prenosu UID je riešená na aplikačnej úrovni. Tieto obmedzenia vytvárajú presný rámec pre interpretáciu výsledkov --- ak sa ukáže zlepšenie odolnosti, ide o~prínos navrhnutých modifikácií v~rámci definovaného prototypu, nie o~všeobecné tvrdenie o~bezpečnosti prístupových systémov.

\begin{table}[!ht]
\renewcommand{\arraystretch}{1.3}
\centering
\caption{Prepojenie úloh zadania s~kapitolami práce}
\label{tab:tasks_mapping}
\small
\begin{tabular}{c p{6.5cm} c}
\toprule
\textbf{Úloha} & \textbf{Obsah} & \textbf{Kapitola} \\
\midrule
1 & Analýza aktuálnych riešení a~srovnávacia základňa & 2 \\
2 & Kryptografické základy, nonce stratégie, HKDF & 2 \\
3 & Návrh a~implementácia prototypu & 3 \\
4 & Experimentálna analýza (7 scenárov, 6 profilov) & 3 \\
5 & Dokumentácia & celá práca \\
\bottomrule
\end{tabular}
\normalsize
\end{table}